% !TeX root = ../nc-quadric-surface-avatar.tex
\section{反標準 avatar の族と二つの十本矢 chart}
\label{sec:qa-avatar-charts}

本節では幾何的な族を直接定義し、その間の開埋め込みを証明する。
代数の同型類から幾何的 avatar への対応は
\cite[定義15.2・定理15.3]{Q12Avatar2026}に記載されているが、
その同型類分類を相対 AS 代数の stack 同値として使用しない。
族上の再埋め込みは\cite[第4--5節]{Q12Weyl2026}との比較も踏まえ、
本稿の補題~\ref{lem:qa-core-curve-vanishing}、
\ref{lem:qa-core-two-pencils}と定理~\ref{thm:qa-core-avatar-square-rotated}で
必要な幾何を自足的に証明し、
開部分の判定に必要なコホモロジー、非被約因子、垂直成分および降下を以下で確認する。
統合ノート\cite[第7.11--7.12項・第9節・第13.3項]{Synthesis2026}は
これらの出典と主張の水準を整理した資料であり、独立の検証の代用とはしない。

\paragraph{AS 側の辞書の対象.}
入力する体上分類の対象は
\cite[第15.1項・定理15.3]{Q12Avatar2026} の $A_Q^{\mathrm{pres}}$ である。
すなわち、局所有限な連結 positively directed $\mathbb Z$-代数で、
大域次元三、Gorenstein parameter 四の AS 正則性を満たし、
\eqref{eq:qa-universal-potential} の十二本矢 potential
（可縮な二次 pair を許す）の展開として表示できるものを、
全整数頂点を固定する代数同型で分類する。
最小化した四型では、右射影 $\mathsf P_i=e_iA$ と右単純 $\mathsf S_i$ に対して
\[
\begin{aligned}
0\longrightarrow\mathsf P_{i-4}
 &\longrightarrow\mathsf P_{i-3}^{\oplus2}\oplus\mathsf P_{i-2}^{\oplus\epsilon_i}\\
 &\longrightarrow\mathsf P_{i-1}^{\oplus2}\oplus\mathsf P_{i-2}^{\oplus\epsilon_i}
 \longrightarrow\mathsf P_i\longrightarrow\mathsf S_i\longrightarrow0
\end{aligned}
\]
という実際に完全な最小分解がある。
$\epsilon_i=0$、奇数 $i$ のみで $\epsilon_i=1$、
偶数 $i$ のみで $\epsilon_i=1$、全 $i$ で $\epsilon_i=1$ が、
順に八本、奇数対角十本、偶数対角十本、十二本に対応する。
右側の AS--Gorenstein 条件は
$\Ext_A^q(\mathsf S_i,\mathsf P_j)=0$ for $(q,j)\ne(3,i-4)$、
$\Ext_A^3(\mathsf S_i,\mathsf P_{i-4})\simeq\kk$ である。
分解の Euler 特性だけを課す条件ではない。
Noether 性はこの定義の追加仮定ではなく、引用する分類の結論である。
potential 表示、trivial pair の係数、周期同型、分解の基底は選択データにしない
\cite[定義2.1・第15.1項]{Q12Avatar2026}\cite[定義48.1]{AvatarVSeven2026}。
この整数頂点固定の同型は、四頂点 graded 代数の同型を射とした
$\mathfrak M_r^{\mathrm{alg}}$ と同じ相対同型層であるとは仮定しない。
本節の幾何的 stack の同型との比較も、体上の同型類辞書と区別する。

\subsection{分類する幾何データ}

基礎体 $\kk$ は標数零の代数閉体とする。
射影束は一次元商を分類する規約を用いる。
次の二つの標識付き曲面を考える。
\[
\begin{array}{c|c|c}
Y&h&f\\\hline
\Sigma _0=\PP^1\times\PP^1&\OO(1,1)&\OO(1,0)\\
\Sigma _2=\PP_{\PP^1}(\OO\oplus\OO(2))&\OO_{\PP(\OO\oplus\OO(2))}(1)&\pi^*\OO_{\PP^1}(1)
\end{array}
\]
いずれも $h^2=2$, $hf=1$, $f^2=0$, $-K_Y=2h$ である。
$\Sigma _2$ の負断面は $S=h-2f$ である。
本節の $\Pic^0$ は Picard 群の単位成分を表す。

\begin{definition}[曲線型の相対 avatar]
\label{def:qa-avatar}
局所 Noetherian な $\kk$-scheme $T$ 上の
$\mathcal A_{12}^{\times}$ の対象を
\[
 (Y\xrightarrow{\pi}B\longrightarrow T;\ h,f;\ D\hookrightarrow Y;\ P)
\]
とする。ここで $B/T$ は $\PP^1$-fibration、$Y/B$ は
$T$ 上 fppf 局所で $\PP_B(\mathcal E)$ と表示される射影直線族であり、
各幾何 fiber の $(Y,h,f)$ は上表のいずれかである。
$h,f$ は相対 Picard の類で、$f$ は指定した ruling に対応する。
$D$ は $T$ 上平坦な相対有効 Cartier 因子で、相対因子類は $2h$ である。
$P$ は $\Pic^0_{D/T}$ の section で、各幾何 fiber 上で非自明とする。
射は ruling、標識、$D$ および $P$ の類を保つ $T$ 上の同型である。
線束の scalar trivialization はデータに加えない。
\end{definition}

この定義は $\Sigma _0$ と $\Sigma _2$ の間で型が跳ぶ族を許す。
実際の線束を用いる式は fppf 局所で解釈し、基底からの線束の tensor は区別しない。
以下で
\[
 p=\pi|_D:D\longrightarrow B,\qquad F=\OO_Y(f)|_D,\qquad H=\OO_Y(h)|_D
\]
と書く。特に $H$ は線束であり、内在的モジュライの記号とは無関係である。
曲面、因子および相対 Picard 類の降下により、定義\ref{def:qa-avatar}は
fppf stack を定める。本節の「開部分 stack」は任意の基底への引戻しが
開部分 scheme となる意味で用いる。

\begin{remark}[有限表示と scalar の扱い]
fppf 局所で $B=\PP^1_T$ とし、$\mathcal E=\pi_*\OO_Y(h)$ を取ると、
各幾何 fiber の分裂型は $\OO(1)^{\oplus2}$ または $\OO\oplus\OO(2)$ である。
評価列は
\[
 0\longrightarrow \mathcal R\otimes\OO_B(-1)
 \longrightarrow \mathcal V\otimes\OO_B
 \longrightarrow\mathcal E\longrightarrow0,
 \qquad\rank\mathcal R=2,\quad\rank\mathcal V=4
\]
となる。したがって線形行列、その余核の射影束上の相対反標準因子、
および相対 Picard 類によって局所的な parameter を記述できる。
この表示の基底変更をすべて新たな自己同型として保持すると、
定義\ref{def:qa-avatar}と異なる stack になる。
本節の射はあくまで標識付き曲面と因子の同型である。
\end{remark}

\subsection{非被約因子を含むコホモロジー}

\begin{lemma}[支持に対する判定]
\label{lem:qa-support}
体上で $D\in|2h|$ とする。このとき
$\omega_D\simeq\OO_D$、$h^0(\OO_D)=h^1(\OO_D)=1$ である。
真の非零部分因子 $0<Z<D$ について $Z(D-Z)\geq2$ が成り立つ。
線束 $N$ の非零切断が正確に $Z\leq D$ に支持されるなら
\[
 \deg_ZN\ \geq\ Z(D-Z).
\]
従って、線束 $L$ がすべての $0<Z\leq D$ について
\begin{equation}
 \deg_ZL+Z(D-Z)>0
 \label{eq:qa-support-vanishing}
\end{equation}
を満たせば $H^1(D,L)=0$ である。
\end{lemma}

\begin{proof}
随伴と $0\to\omega_Y\to\OO_Y\to\OO_D\to0$ により
双対化層と構造層のコホモロジーを得る。
$\Sigma _0$ では $D=2f+2g$、$Z=af+bg$ と書けて
\[
 Z(D-Z)=2(a+b-ab),\qquad 0\leq a,b\leq2.
\]
$Z=0,D$ を除いたすべての場合で右辺は二以上である。
$\Sigma _2$ では $D=2S+4f$、$Z=aS+bf$ と書くと
\[
 Z(D-Z)=2a^2+2b(1-a),\qquad0\leq a\leq2,\quad0\leq b\leq4.
\]
$a=0$ では $b\geq1$、$a=1$ では右辺が二、
$a=2$ では真部分因子なら $b\leq3$ なので、同じ結論を得る。

切断の generic annihilator によって、その正確な支持である部分因子 $Z$ を取る。
各既約成分の generic な Artin 局所環で切断の消滅次数を取り除いて $Z$ の重複度を定めると、
切断は $N|_Z(-(D-Z))$ の、各 generic point で可逆な切断になる。
$Z$ は Cohen--Macaulay なので、この切断が定める
$\OO_Z\to N|_Z(-(D-Z))$ は単射で、余核は有限長である。
Euler 特性を比較すれば不等式を得る。
これは部分因子を重複度込みで取る議論であり、$D$ の被約性を使わない。
最後は $\omega_D\simeq\OO_D$ による Serre 双対を用い、
$L^{-1}$ の切断にこの不等式を適用すればよい。
\end{proof}

\begin{lemma}[再埋め込みに必要な消滅]
\label{lem:qa-vanishing}
$R\in\Pic^0(D)$ とする。このとき
\[
 H^1(D,HR)=H^1(D,HF^{-1}R)=0,\qquad
 \chi(HR)=4,\quad\chi(HF^{-1}R)=2.
\]
また $R\not\simeq\OO_D$ なら $H^*(D,R)=0$ である。
\end{lemma}

\begin{proof}
$R$ は各部分因子上で次数零である。
$h$ は nef で全次数が四なので、補題\ref{lem:qa-support}を適用できる。
$\Sigma _0$ では $h-f$ も nef で全次数が二である。
$\Sigma _2$ では $Z=aS+bf$ 上で
\[
 \deg_Z(h-f)+Z(D-Z)=2a^2-a+3b-2ab.
\]
$a=0,1,2$ に対する右辺はそれぞれ $3b$, $1+b$, $6-b$ であり、
$Z\neq0$ の許される範囲で正である。Riemann--Roch が Euler 特性を与える。
$R$ の非零切断は真部分因子には支持されない。
全体に支持される切断の余核は次数零なので零であり、$R$ は自明になる。
$R^{-1}$ にも同じ議論を適用し、Serre 双対を使う。
\end{proof}

\subsection{Picard twist による相対的な再埋め込み}

\begin{proposition}[垂直成分がある場合も含む直像]
\label{prop:qa-pushforward}
定義\ref{def:qa-avatar}の族で、$R\in\Pic^0_{D/T}(T)$ を取る。
fppf 局所で線束によって表示すると
\[
 Rp_*(HR)\simeq\mathcal E_R:=p_*(HR)
\]
は $B$ 上の階数二の局所自由層であり、任意の基底変換と両立する。
各幾何 fiber $B_t\simeq\PP^1$ 上で次数二かつ大域生成である。
評価写像は閉埋め込み
\[
 j_R:D\hookrightarrow Y_R:=\PP_B(\mathcal E_R)
\]
を定め、像は相対反標準 Cartier 因子である。
新しい標識の制限は $F_R=F$、$H_R=HR$ である。
\end{proposition}

\begin{proof}
まず体上で $b\in B$、$G_b=\pi^{-1}(b)$ とする。
$G_b$ が $D$ の成分でなければ $D\cap G_b$ は長さ二の scheme であり、
導来 fiber の切断は二次元で次数零に集中する。
$G_b\leq D$ なら、局所 parameter $t$ と $D$ の方程式 $t^m a$、$t\nmid a$ を用いると
\[
 \operatorname{ann}_{\OO_Y/(t^ma)}(t)=(t^{m-1}a)/(t^ma).
\]
したがって $L=HR$ の $B$ 方向の導来制限は
\[
 \mathcal H^{-1}=L|_{G_b}\otimes\OO_{G_b}(-D+G_b)\simeq\OO_{\PP^1}(-1),
 \qquad
 \mathcal H^0=L|_{G_b}\simeq\OO_{\PP^1}(1)
\]
を持つ。ここで $\deg_{G_b}L=1$、$D G_b=2$、$G_b^2=0$ を使った。
次数 $-1$ の層は acyclic であるから、この場合も直像の導来 fiber は $\kk(b)^2[0]$ となる。
垂直成分を持つ $p$ は有限でも平坦でもないが、この計算にはその仮定を要しない。

$i_*L$ は Cartier immersion $i:D\hookrightarrow Y$ の線束の直像なので perfect である。
滑らかな射影的射 $\pi$ による perfect pushforward と導来基底変換を用いると、
$R\pi_*i_*L$ は perfect である。
すべての導来 fiber が同じ階数二の次数零複体なので、局所的に可逆 minor で
可縮項を消去することにより、階数二の局所自由層となる。
族の場合、$i_*L$ は $T$ 上平坦であり、まず $T$ の幾何 fiber、次に $B_t$ の一点へ
制限すれば同じ計算になる。このため非被約な基底も含めて結論が成り立つ。

補題\ref{lem:qa-vanishing}と Riemann--Roch により
\[
 \deg\mathcal E_R=\chi(HR)-2=2,\qquad
 H^1(B,\mathcal E_R(-1))=H^1(D,HF^{-1}R)=0.
\]
$\PP^1$ 上の分裂定理より、分裂型は $\OO(1)^{\oplus2}$ または $\OO\oplus\OO(2)$ である。

各 $B$-fiber での評価は、長さ二の scheme 上の完全切断の評価、
または $\OO_{\PP^1}(1)$ の評価である。従って全射で、誘導される fiber の射は閉埋め込みである。
Nakayama の補題から全評価も全射となる。$j_R$ は proper かつ quasi-finite なので有限であり、
有限射の基底変換と再度の Nakayama により
$\OO_{Y_R}\to(j_R)_*\OO_D$ は全射である。
体上では像は滑らかな曲面内の Cohen--Macaulay な余次元一部分 scheme なので Cartier である。
像の類を $a h_R+b f_R$ とすると、ruling との交点から $a=2$、
$h_R$ との交点 $\deg_D(HR)=4$ から $b=0$ を得る。
族では $D/T$ の平坦性と各 fiber の Cartier 性によって相対 Cartier 性が従う。
各 fiber で自明な差の線束は、$H^1(\OO_{Y_t})=0$ と基底変換により基底から来るので、
相対 Picard 類も $2h_R$ である。
\end{proof}

\begin{theorem}[族上の幾何的並進]
\label{thm:qa-translation}
前命題において $R=P^n$ とした構成は $\mathcal A_{12}^{\times}$ の自己同値
$\mathsf T^n$ を定める。自然な同型
\[
 \mathsf T^m\mathsf T^n\simeq\mathsf T^{m+n},\qquad
 \mathsf T^0\simeq\operatorname{id}
\]
があり、逆操作は $P^{-n}$ を用いる同じ構成である。
\end{theorem}

\begin{proof}
抽象的な $D$、ruling $p:D\to B$ および $P$ は変更せず、$H$ を $HP^n$ に変える。
二回適用した直像は $p_*(HP^nP^m)$ なので、tensor の結合則が合成則を与える。
$R=\OO_D$ の場合、制限列
\[
 0\longrightarrow\OO_Y(-h)\longrightarrow\OO_Y(h)\longrightarrow i_*H\longrightarrow0
\]
と $R\pi_*\OO_Y(-h)=0$ が元の $\mathcal E=\pi_*\OO_Y(h)$ を回収する。
相対的には左項に基底線束の tensor が入ることがあるが、同じ消滅が成り立つ。
評価の自然性により射影束だけでなく $D$ の埋め込みも回収する。
線束の代表を基底線束で tensor すると直像も同じ線束で tensor されるため、
射影束および相対標識は変わらない。
$q:D\to T$ について、各幾何 fiber で $h^0(\OO_D)=1$ であるから、
コホモロジーと基底変換により $q_*\OO_D=\OO_T$ が任意の基底変換後も成立する。
従って線束代表の自己同型は基底からの scalar であり、射影化で消える。
従って構成と合成同型は fppf 降下し、射と基底変換に整合する。
\end{proof}

この定理は\cite[定理5.4]{Q12Weyl2026}の操作である。
ここで証明したのは幾何的 avatar の stack 上の並進であり、
無限二面体群のすべての関係を stack 上へ持ち上げる主張ではない。

\subsection{二つの開 chart と四つの層}

\begin{definition}[幾何的な八本矢・十本矢 chart]
\label{def:qa-open-charts}
$\mathcal U\subset\mathcal A_{12}^{\times}$ を、全幾何 fiber が $\Sigma _0$ である部分とする。
\[
 \mathcal A_{10,\mathrm e}^{\times}:=\mathcal U,
 \qquad
 \mathcal A_{10,\mathrm o}^{\times}:=\mathsf T^{-1}(\mathcal U),
 \qquad
 \mathcal A_8^{\times}:=\mathcal U\cap\mathsf T^{-1}(\mathcal U)
\]
と置く。添字 $\mathrm e,\mathrm o$ は、後述する標準表示で残る偶数対角・奇数対角に対応する。
ここでは定義自体を幾何データだけで与えている。
\end{definition}

\begin{theorem}[avatar の開部分 stack]
\label{thm:qa-open-charts}
次の射は開埋め込みである。
\[
 \mathcal A_8^{\times}\hookrightarrow
 \mathcal A_{10,\mathrm o}^{\times}\hookrightarrow\mathcal A_{12}^{\times},
 \qquad
 \mathcal A_8^{\times}\hookrightarrow
 \mathcal A_{10,\mathrm e}^{\times}\hookrightarrow\mathcal A_{12}^{\times}.
\]
幾何点での層は次の表で与えられる。
\begin{center}
\begin{tabular}{c|c|c}
曲面&曲線型の条件&属する層\\\hline
$\Sigma _0$&$P\not\simeq F^2H^{-1}$&$\mathcal A_8^{\times}$\\
$\Sigma _0$&$P\simeq F^2H^{-1}$&$\mathcal A_{10,\mathrm e}^{\times}\setminus\mathcal A_8^{\times}$\\
$\Sigma _2$&$D\cap S=\varnothing$&$\mathcal A_{10,\mathrm o}^{\times}\setminus\mathcal A_8^{\times}$\\
$\Sigma _2$&$S\leq D$&$\mathcal A_{12}^{\times}\setminus
(\mathcal A_{10,\mathrm e}^{\times}\cup\mathcal A_{10,\mathrm o}^{\times})$
\end{tabular}
\end{center}
最後の行は幾何点集合として閉 locus である。
全 fiber がこの型となる族の充満部分関手に、閉埋め込みを自動的に与える主張ではない。
\end{theorem}

\begin{proof}
$\mathcal E=\pi_*\OO_Y(h)$ の二つの分裂型は
\[
 h^0(B_t,\mathcal E_t(-2))=
 \begin{cases}0&Y_t=\Sigma _0,\\1&Y_t=\Sigma _2\end{cases}
\]
によって区別される。コホモロジーの上半連続性から $\mathcal U$ は開である。
自己同値 $\mathsf T$ による逆像と交わりも開なので、すべての開埋め込みを得る。
残るのは幾何点での条件である。

$\Sigma _0$ で $H=FG$ と書くと
\[
 h^0(B,\mathcal E_P(-2))=h^0(D,GF^{-1}P).
\]
$Z=af+bg$ に対し、$GF^{-1}P$ の次数は $a-b$ である。
真部分因子では
\[
 Z(D-Z)-\deg_Z(GF^{-1}P)=a+3b-2ab>0.
\]
よって補題\ref{lem:qa-support}により、この線束の非零切断は全体に支持される。
全次数は零なので、非零切断があることは線束が自明であることと同値である。
従って $\mathsf T(Y)=\Sigma _2$ となる条件はちょうど $P\simeq F^2H^{-1}$ である。

$\Sigma _2$ で $D\cap S=\varnothing$ なら $H\simeq F^2$ だから
$h^0(D,HF^{-2}P)=h^0(D,P)=0$ であり、$\mathsf T(Y)=\Sigma _0$ となる。
$S\leq D$ なら $C=D-S$ と置く。$C\sim h+2f$ であり、
$\OO_Y(-C)$ は $\pi$ に沿って次数 $-1$ なので acyclic である。
従って $h^0(\OO_C)=1$, $h^1(\OO_C)=0$、標数零より $\Pic^0(C)=0$ を得る。
残余因子の完全列は
\[
 0\longrightarrow P|_C\longrightarrow\OO_D(S)P
 \longrightarrow\OO_S(-2)\longrightarrow0.
\]
$P|_C\simeq\OO_C$ なので $h^0(D,HF^{-2}P)=1$ である。
従って $\mathsf T(Y)=\Sigma _2$ のままである。
最後に $D S=0$ であるから、有効因子 $D$ は $S$ を含むか $S$ を避けるかのいずれかである。
これですべての場合を得る。閉性は二つの開部分の和の補集合であることによる。
\end{proof}

\begin{remark}[矢数の名称と体上の辞書]
幾何的な開埋め込みと表の分解は、上の証明の通り
射影束の分裂型、Picard 類および垂直成分条件から従う。
各層を $8,10,10,12$ 本の最小矢数と対応させる体上の解釈は、
$A_Q^{\mathrm{pres}}$ を対象とする
\cite[定理15.3]{Q12Avatar2026}の同型類分類を用いる。
同じ矢数の相対 QP 表示を比較する際の階数条件は、
上の幾何学的な層の名称からは
推論せず、各代数表示の積と関係について検証する。
\end{remark}


\begin{remark}[「十本矢全体」の二通りの意味]
最小十本矢の二つの層の和は $\mathcal A_8^{\times}$ を含まない。
これに対し上の $\mathcal A_{10,\mathrm o}^{\times}$ と
$\mathcal A_{10,\mathrm e}^{\times}$ は、可縮な pair を許す十本矢表示の
幾何的 chart に対応し、八本矢部分を開部分として含む。
二つの十本矢 chart の和は、最深の十二本矢層の補集合である。
この区別なしに「矢の少ないモジュライが矢の多いモジュライの開部分である」
と述べることはできない。
\end{remark}

\subsection{直積型 collection と全 square avatar}

$\mathcal A_{\mathrm{sq}}^{\times}$ を、順序付きの二つの ruling を持つ
$\Sigma _0$-族、その相対 $(2,2)$ Cartier 因子 $D$、および fiberwise 非自明な
$P\in\Pic^0_{D/T}$ の stack とする。
曲線上の ruling の制限を $F_s,G_s$ と書く。
これは $h=f+g$ と置くことによって $\mathcal U$ と同じ幾何的 stack である。

\begin{corollary}[全 square chart の相対的な対応]
\label{cor:qa-square-chart}
幾何的 stack の同値
\[
 \mathcal A_{\mathrm{sq}}^{\times}
 \xrightarrow[\ \mathsf T^{-1}\ ]{\ \sim\ }
 \mathcal A_{10,\mathrm o}^{\times}
\]
があり、制限線束は
\begin{equation}
 F=F_s,\qquad H=F_sG_sP^{-1},\qquad P=P
 \label{eq:qa-square-dictionary}
\end{equation}
となる。逆には $\mathsf T$ によって得られる $\Sigma _0$ 上で
\[
 F_s=F,\qquad G_s=HPF^{-1}
\]
を取る。square 座標で八本矢部分となる条件は
$P\not\simeq G_sF_s^{-1}$ であり、壁 $P\simeq G_sF_s^{-1}$ は
負断面を避ける $\Sigma _2$ 上の奇数対角十本矢層に移る。
\end{corollary}

\begin{proof}
定義から $\mathsf T^{-1}:\mathcal U\to\mathsf T^{-1}(\mathcal U)$ は同値であり、
その線束の式が\eqref{eq:qa-square-dictionary}である。
逆の式は $H_s=HP=F_sG_s$ による。
元の $\Sigma _0$ で $\mathsf T^{-1}$ が $\Sigma _2$ を与える条件は、
定理\ref{thm:qa-open-charts}の計算を $P^{-1}$ に対して行って
$P^{-1}=F_sG_s^{-1}$、すなわち $P=G_sF_s^{-1}$ である。
その像は $\mathcal A_{10,\mathrm o}^{\times}$ にあるので負断面を避ける。
\end{proof}

この系では反標準因子を滑らか、被約または既約と仮定していない。
従って「十本でちょうど全ての quadric の divisor と $\Pic^0$ が出る」
という記述は、\emph{奇数対角の非最小十本矢 chart と全 square avatar の対応}
として正確に定式化できる。最小十本矢だけの主張ではなく、$P=\OO_D$ も除いている。

代数側の square foundation は
\[
 B_{\mathrm{sq}}=\End_{\to}(P,F_s,G_s,F_sG_s)
\]
であり、可換な基準 collection は
$(\OO,\OO(1,0),\OO(0,1),\OO(1,1))$ である
\cite[定義1.2・定理1.3]{AvatarVSeven2026}。
最初の二項の左 mutation は曲線上で
\[
 (P^2F_s^{-1},P,G_s,F_sG_s)
\]
を与える。共通に $F_sP^{-1}$ を tensor すると
\[
 (P,F_s,F_sG_sP^{-1},F_s^2G_sP^{-1})=(P,F,H,HF)
\]
となり、上の幾何的辞書と一致する。
mutationによる指定型AS代数の対応は定理~
\ref{thm:qa-as10-global-bijection}と
\ref{thm:qa-as10-relative-families}で証明する。
元の体上および仮定付き相対対応は
\cite[定理51.4・系51.5、定理53.2]{AvatarVSeven2026}にある。
幾何的な開埋め込みは独立に証明した。

\subsection{滑らかな楕円型における三者の明示的辞書}
\label{subsec:qa-square-elliptic-pointed-dictionary}

\begin{theorem}[VdB楕円四つ組とsquare avatarの同型類の全単射]
\label{thm:qa-elliptic-vdb-square-avatar-bijection}
$k$を標数零の代数閉体とする。
滑らかな楕円曲線$E$上のVan den Bergh許容四つ組
$(E,L_0,L_1,L_2)$で、各$L_i$は次数二、
\[
 L_1L_0^{-1}\ne\OO_E,\qquad
 L_2L_0^{-1}\ne\OO_E,\qquad
 L_1L_2^{-1}\ne\OO_E
 \tag{VdB-adm}\label{eq:qa-vdb-square-admissibility}
\]
を満たすものの、曲線・三線束を保つ同型類を考える。
これは、\emph{順序付き}二rulingを持つ
滑らかなsquare avatar
$(Y=\PP^1\times\PP^1,D,P)$で
\[
 D\in|\OO_Y(2,2)|,\qquad D\text{ は滑らか},\qquad
 P\in\Pic^0(D)\setminus\{\OO_D\},\qquad
 P\not\simeq G_sF_s^{-1}
 \tag{Sq-cub}\label{eq:qa-square-elliptic-cubic-open}
\]
を満たすものの同型類と自然に全単射する。
ここで$F_s=\OO_D(1,0)$、
$G_s=\OO_D(0,1)$とし、
二rulingの順序は同型の中で保持する。

全単射の\emph{両方向の式}は
\[
 \boxed{
 \begin{aligned}
 (E;L_0,L_1,L_2)&\longmapsto
  (D\hookrightarrow\PP H^0(F_s)^\vee
           \times\PP H^0(G_s)^\vee;\ P),\\
 F_s&=L_2,\qquad
 G_s=L_1L_2L_0^{-1},\qquad
 P=L_2L_0^{-1},\\
 (D;F_s,G_s,P)&\longmapsto
  (E=D;\ L_0=F_sP^{-1},\
   L_1=G_sP^{-1},\ L_2=F_s).
 \end{aligned}}
 \tag{Sq-VdB}\label{eq:qa-square-vdb-explicit}
\]
で与えられる。
この同型類の比較は、周期同型や根を
追加標識にしないVdBのcubic$\mathbb Z$-代数
の幾何的分類を通じて、
\emph{選んだcubic normal foundation}の
相対的な$k$点の辞書になる。

差パラメータ
$\eta=L_1L_0^{-1}$、
$\delta=L_2L_0^{-1}$を用いれば、
\[
 \eta=G_sF_s^{-1},\qquad \delta=P.
 \tag{Sq-diff}\label{eq:qa-square-vdb-differences}
\]
\end{theorem}
\begin{proof}
まず$L_0,L_1,L_2$から
$P=L_2L_0^{-1}$、
$F_s=L_2$、
$G_s=L_1L_2L_0^{-1}$を作ると、
$F_s,G_s$はいずれも次数二で、
$G_sF_s^{-1}=L_1L_0^{-1}$である。
式~\eqref{eq:qa-vdb-square-admissibility}より
$F_s\not\simeq G_s$である。
滑らかな楕円曲線上の相異なる次数二線束の
完全線形系は、それぞれ二重被覆
$E\to\PP^1$を与え、その積は
$\PP^1\times\PP^1$への閉埋込みである。
像は$(2,2)$因子で、反標準因子となる。
また$P\ne\OO_E$、
$P\not\simeq G_sF_s^{-1}$である。
これは曲線同型を保ち、線形系の基底の選択を
$\PGL_2^2$で割ればavatarの同型類を定める。

逆にsquare avatarから
$L_0=F_sP^{-1}$、
$L_1=G_sP^{-1}$、
$L_2=F_s$を作ると、
式~\eqref{eq:qa-square-elliptic-cubic-open}により
$\eta=G_sF_s^{-1}\ne\OO_E$、
$\delta=P\ne\OO_E$、
$\eta\ne\delta$となる。
従ってVdBの許容条件が成立する。
二つの式を順次適用すると三線束と
二つのrulingの完全線形系をそのまま回収するので、
互いに逆である。
VdBの幾何的quintuple復元
\cite[Theorem 5.5.10]{VdB2011}を適用すれば
最後の解釈も従う。
\end{proof}

\begin{theorem}[O-markingを保つruling反射と楕円型の商全単射]
\label{thm:qa-elliptic-avatar-pointed-quotient}
前定理のsquare avatarについてさらに
\[
 G_sF_s^{-1}\not\simeq P^2
 \tag{two-cub}\label{eq:qa-square-two-cubic-open}
\]
を課した開を$\mathcal A_{\mathrm{sq,ell}}^{2c}$とする。
この開上では
\[
 \boxed{\quad
 r_1(F_s,G_s,P)
  =(G_sP^{-1},F_sP,P)
 \quad}
 \tag{Sq-R}\label{eq:qa-square-pointed-reflection}
\]
は、右辺の二つの次数二線束が定める
新しい$(2,2)$埋込みを用いて、
幾何的avatarの同型類の対合を定める。
前定理の式~\eqref{eq:qa-square-vdb-explicit}の下で、
これは論文Bの
\cite[\texttt{lem:elliptic-two-reflections}]{QuadricModuli2026}
の\emph{構造対象$\OO$を保つ}ruling交換
\[
 (L_0,L_1,L_2)\longmapsto
 (L_0L_1L_2^{-1},L_2,L_1)
\]
に一致する。
従って、$y=(\mathcal C,\OO)$が滑らかな
楕円型かつ二つのnormal foundationがともに
cubic型となる\emph{幾何的なO-marked同型類集合}は
\[
 \boxed{\quad
 \bigl|\mathcal A_{\mathrm{sq,ell}}^{2c}(k)\bigr|
    /\langle r_1\rangle
 \ \xrightarrow{\ \sim\ }\
 \bigl|\mathfrak H_{2,\mathrm{ell}}^{2c}(k)\bigr|
 \quad}
 \tag{Sq-H2}\label{eq:qa-square-pointed-bijection}
\]
と全単射する。
左辺はavatar同型類の集合をとってから
$r_1$の軌道で割ることを意味する。
ここでは$S_2$の商の全相対stack同値を
この幾何的全単射だけから推論しない。

一方、helixの一段の移動に対応する
\[
 r_0(F_s,G_s,P)=(G_s,F_sP,P)
 \tag{Sq-T}\label{eq:qa-square-unpointed-shift}
\]
はVdB四つ組の
$(L_0,L_1,L_2)\mapsto(L_1,L_2,L_3)$に
対応するが、一般に$\OO$を変更する。
従って\emph{O-marked}同型類を得るために
$r_0$まで同一視してはならない。

この商の記述は、一般の二つのavatarの
$r_1$-軌道が二元であることを要求しない。
特殊な二捩れ慣性点では二つが同型になる。
\end{theorem}
\begin{proof}
式~\eqref{eq:qa-square-vdb-explicit}を
$r_1$の右辺へ適用すると
\[
 \begin{aligned}
 L'_0&=G_sP^{-2}
       =L_0L_1L_2^{-1},\\
 L'_1&=F_s=L_2,\qquad
 L'_2=G_sP^{-1}=L_1
 \end{aligned}
\]
を得る。新しい二ruling線束が相異なる条件は
$G_sP^{-1}\not\simeq F_sP$、
すなわち式~\eqref{eq:qa-square-two-cubic-open}
に他ならない。
もう一度適用すると$(F_s,G_s,P)$を回収するため
対合である。

論文Bの二つの内在的foundationの系
\cite[\texttt{cor:two-intrinsic-foundations}]{QuadricModuli2026}と、
pointed groupoid比較の系
\cite[\texttt{cor:h2-point-groupoid-foundation}]{QuadricModuli2026}は、指定$\OO$から始まるnormal foundationは
exact同値の下でちょうど二択であり、
二つともcubic型なら他方は上記$R$で得られること、
またfoundation同型はpointed exact同型を
完全に検出することを示す。
Van den Berghの四つ組復元を合わせれば、
二cubic楕円型のpointed同型類に関する
核と像はちょうど$r_1$の同型類軌道となる。
これが式~\eqref{eq:qa-square-pointed-bijection}である。

同様に$r_0$の右辺から得る
$(L'_0,L'_1,L'_2)$は
\[
 (G_sP^{-1},F_s,G_s)=(L_1,L_2,L_3),
 \quad L_3=L_0^{-1}L_1L_2
\]
である。論文Bの同じ補題では、この添字移動は
underlying Abel圏の同値類は保つが、
一般に構造対象$\OO$を保たない。
従って$r_0$を割るのは別の
unpointed moduli問題である。
\end{proof}

\begin{corollary}[OU potentialからsquare avatar、O-marked対象への一般次数]
\label{cor:qa-ou-potential-square-avatar-pointed-degrees}
論文C
~\ref{thm:qa-ou-rigid-potential-quintuple-generic-etale}
の正則半単純な滑らかな楕円開を、
前定理の二cubic開~\eqref{eq:qa-square-two-cubic-open}
および一般慣性同型の稠密開へ制限する。
この開ではpotentialの幾何的四重表示が
VdB四つ組$(E,L_0,L_1,L_2)$を与え、
式~\eqref{eq:qa-square-vdb-explicit}により
\emph{square avatarへの相対的な対応}がある。
全ての線束・曲線が相対的に構成される
étale chart上では、これは基底変更と可換な
代数スタックの射となり、
相対的な曲線同型を通じて降下する。
一方、対応の幾何的fiberの一般次数は
\[
 \boxed{
 \begin{array}{ccc}
 \text{有効potential表示}
   &\xrightarrow{\;4\;}&
  \text{順序付きsquare avatar}\\
 &&\big\downarrow\vcenter{\rlap{$\scriptstyle 2$}}\\
 &&\text{O-marked }\mathfrak H_2
 \end{array}}
 \tag{Pot-Sq-H2}\label{eq:qa-potential-square-pointed-degrees}
\]
であり、potentialからO-marked対象への
一般次数は$4\cdot2=8$である。
ここで上の$4$はpotential/quintupleの
\emph{表示間}の一般四重射、
下の$2$は二つのnormal foundationの忘却であり、
$\Pic^0$の四つの平方根の選択を
追加標識として数えた次数ではない。
\end{corollary}
\begin{proof}
Cの
~\ref{thm:qa-ou-rigid-potential-quintuple-generic-etale}
による一般有限étale四重射は、
指定された幾何quintupleの相対四重テンソルを与える。
滑らかな正則半単純開では、C
~\ref{cor:qa-ou-relative-regular-semisimple-potential}
の点曲線と二rulingの相対構成が、
基底変更と可換な三つの次数二線束を与える。
それらのテンソル積と二つの完全線形系で
式~\eqref{eq:qa-square-vdb-explicit}の
square avatarを相対的に作る。
曲線・線束の同型を保つので基底変更と降下に
関して自然である。
VdBの四つ組の幾何的同型類とsquare avatarの
全単射は定理
~\ref{thm:qa-elliptic-vdb-square-avatar-bijection}で証明した。
これとCの
~\ref{cor:qa-quintuple-h2-generic-degree-two}の
二つのfoundationの一般次数を合成すれば
表示された幾何一般次数が従う。
\end{proof}

\begin{remark}[一cubic境界と古典点での辞書の射程]
\label{rem:qa-square-avatar-boundary-cautions}
$G_sF_s^{-1}=P^2$では、$r_1$後の
$F'_s=G_sP^{-1}$と$G'_s=F_sP$が一致し、
square$(2,2)$埋込みを構成できない。
これは論文C
~\ref{cor:qa-ou-one-cubic-fiber-and-nonfiniteness}
の一cubic境界であり、
二つ目のnormal foundationはnoncubicとして
内在的な選択被覆
~\ref{thm:qa-global-two-ruling-etale-cover}
に残る。

また$P=\OO_D$を任意の$(2,2)$因子に
そのまま許すと、これらの曲線から得る
可換な四項foundation代数は$D$に依存せず、
同じ可換quadricを重複して数える。
従って全同型類の全単射には古典点を
別途一点として追加し、元のavatarの
$P=0$座を単に追加するだけでは足りない。
$\Sigma_2$の古典型は別のweighted
stackを生み、無条件には$\HHstack$に入らない。
$\Sigma_2$曲線型でのO-marked適合性、
非被約反標準因子を含む全対応、
およびpotentialの安定化・自己同型群の比較は
式~\eqref{eq:qa-square-pointed-bijection}
の証明には含まれない。
\end{remark}

\begin{proposition}[二つのwallは異なり、noncubic avatar壁には追加認識が必要]
\label{prop:qa-avatar-walls-and-noncubic-dimension}
滑らかな楕円型square avatarについて、
$\eta=G_sF_s^{-1}$、$\delta=P$と書く。
次の二つのwallを区別する。
\begin{enumerate}
\item $\eta=2\delta$、$\delta\notin E[2]$：
式~\eqref{eq:qa-square-vdb-explicit}の四つ組は
geometric cubic型であるが、
$r_1$後の第二foundationはnoncubicになる。
したがって元のO-marked点は
\emph{cubic locusに属する}。
\item $\eta=\delta\ne0$：
式~\eqref{eq:qa-square-vdb-explicit}の
$L_1=L_2$なので元のfoundationが
geometric cubic型ではない。
このwallはsquareの八本矢開の補集合であり、
曲線型avatarの十本矢・$\Sigma_2$ chartへ
移行する。
\end{enumerate}
後者の全smooth楕円avatar同型類の
パラメータ空間は、楕円曲線の$j$と
非零点$\delta\in E$を局所座標として
\emph{二次元}である。
一方、論文Bの内在noncubic境界
$\mathfrak N\subset\HHstack$は
\emph{純一次元}のArtinスタックである。
従ってこのwall\emph{全体}を
内在noncubic境界と表現可能quasi-finiteな
モジュライ同型で結び付けることはできない。
O-marked適合性はこのwall上で
さらに制限されるか、あるいは
同じpointed対象へ正の相対次元を持って
潰れるかのいずれかである。

特に非零二捩れ$P$の条件$P^2=\OO_E$は
楕円曲線のmoduli上で一次元を定め、
次に検討すべき\emph{候補}ではあるが、
それが内在noncubic境界をちょうど尽くすことは
ここでは証明していない。
\end{proposition}
\begin{proof}
最初のwallでは$L_0=F_sP^{-1}$、
$L_1=G_sP^{-1}$、$L_2=F_s$が
互いに相異なり、VdB許容条件を満たす。
$r_1$による新しい二線束は
$G_sP^{-1}=F_sP$と$F_sP$で一致するので、
第二はnoncubicとなる。
これはC
~\ref{cor:qa-ou-one-cubic-fiber-and-nonfiniteness}
の一cubic境界に一致する。
次のwallでは$G_s=F_sP$だから
$L_1=G_sP^{-1}=F_s=L_2$となる。
$\mathcal A_8^\times$の条件
~\eqref{eq:qa-square-elliptic-cubic-open}
に反することも直接に分かる。
$\Sigma_2$との対応は
系~\ref{cor:qa-square-chart}に従う。

最後の次元を計算する。
楕円曲線$E$のmoduliは一次元で、
$F_s\in\Pic^2(E)$の選択は曲線の
translationによって動かせる。
標数零では$[2]:E\to E$が全射なので、
このtranslationの軌道は$\Pic^2(E)$全体である。
$\eta=\delta$の下では$G_s=F_sP$だから
残るパラメータは
$E$と$0\ne P\in\Pic^0(E)$で、
そのmoduliの一般次元は$1+1=2$。
一方、論文B
\cite[\texttt{thm:noncubic-boundary-smooth-codim-two}]{QuadricModuli2026}
は全$\HHstack$が純次元三で
$\mathfrak N$が滑らかな余次元二、
すなわち純次元一であることを証明している。
有限型quasi-finite射の像は源の各成分と同じ
次元を持つので、この二次元wall全体を
一次元の境界にquasi-finiteに写すことはできない。
最後の二捩れ条件の次元も
有限étaleな$E[2]\setminus0$から従う。
\end{proof}

\begin{theorem}[二捩れsquare壁とnoncubic二根対の幾何同型類の全単射]
\label{thm:qa-avatar-two-torsion-root-pair-bijection}
$k$を標数零の代数閉体とする。
次の二集合を考える。
\begin{enumerate}
\item 滑らかな楕円曲線$E$、次数二線束$F_s$、
非自明な$P\in\Pic^0(E)[2]$の三つ組
$(E,F_s,P)$の曲線・線束・Picard類を
保つ同型類集合。$G_s=F_sP$とすると
$F_s\not\simeq G_s$だから、二完全線形系は
$(2,2)$埋込みを定め、対応するsquare avatarは
壁$P=G_sF_s^{-1}$上にある。
\item 射影直線$B\simeq\PP^1$上の、
互いに交わらない\emph{被約}二点因子の
非順序対$\{B_1,B_2\}$で、
$\deg B_i=2$かつ$B_1+B_2$は
四つの相異なる点からなるものの、
$\PGL_2$同型類集合。
\end{enumerate}
この二集合の間に自然な全単射がある。
対応は、$F_s$の二重被覆
\[
 \pi_F:E\longrightarrow\PP H^0(E,F_s)^\vee
\]
の四つの分岐点に対し、
$P$に対応する非零二捩れ平行移動が
誘導する二組の分割を取ることで与えられる。
逆方向には、四分岐点$B_1+B_2$の
二重被覆$\pi_F:E\to B$を取り、
その分割に対応する唯一の
非零$E[2]$の元を$P$とする。
このとき$F_s=\pi_F^*\OO_B(1)$である。

論文B
\cite[\texttt{thm:noncubic-intrinsic-presentation-equivalence}]{QuadricModuli2026}
の非cubic根対表示によれば、
従ってこの二捩れsquare壁の幾何的avatar同型類は、
\emph{四根が相異なるnoncubic O-marked H2境界}
の同型類とも全単射する。

最後の全単射は\emph{論文Bの根対分類を介した}
同型類の対応である。
この幾何辞書だけから
avatar側のAS$\mathbb Z$-代数の
$\qgr$とそのnoncubic O-marked heartの
同値、または対応するモジュライstackの射は
まだ結論しない。
\end{theorem}
\begin{proof}
次数二の線束$F_s$は楕円曲線から$\PP^1$への
二重被覆を与える。
その被覆対合を$\iota_F$と書くと、
その固定点集合$R_F$は四点からなり、
像$B_F$が四つの相異なる分岐点となる。
$E$は標識なしでも$\Pic^1(E)$のtorsorで、
$P\in\Pic^0(E)[2]\setminus0$は
$E$上の固定点のない位数二平行移動$t_P$を
定める。
次数二線束に対する平行移動の作用は
$t_P^*F_s\simeq F_sP^{\pm2}\simeq F_s$なので、
$t_P$は$\iota_F$と可換であり、
$R_F$を保つ。
その四点への作用は自由であり、
ちょうど二つの長さ二の軌道に分ける。
二つの軌道の$\pi_F$による像が
$B_F=B_1+B_2$の二分割となる。
これは線束の基底・楕円原点の選択に依存せず、
曲線と線束の同型に自然である。

逆に四つの相異なる分岐点$B_1+B_2$から、
$B$上の二重被覆は、
$\mathcal O_B(2)$とその四点を零点に持つ
切断を用いて、被覆同型を除き一意に構成される。
$F_s=\pi_F^*\OO_B(1)$は次数二である。
四つの分岐点は$E$上の被覆対合の固定点と
一対一に対応する。任意の一つを楕円原点に
選べば、その四点が$E[2]$となる。
$E[2]\simeq(\mathbb Z/2)^2$の各非零元の
平行移動は四点を二つの組に分割し、
三つの非零元と三つの二分割との間に
自然な全単射を与える。
従って指定の$\{B_1,B_2\}$から
一意の非零二捩れ$P$を回収する。
原点を変えても$\Pic^0(E)$における$P$の
同型類は変わらない。
二方向の構成が互いに逆であることも
この四点の分割から直ちに従う。

具体的に、適当な射影座標で
$B_1=\{0,\infty\}$、$B_2=\{1,\lambda\}$、
$\lambda\notin\{0,1,\infty\}$とすれば、
$E$はLegendre曲線
\[
 y^2=x(x-1)(x-\lambda)
\]
で与えられる。$B_1,B_2$をそれぞれ
集合として保ち、その二点を交換する
$\PP^1$の対合は$x\mapsto\lambda/x$である。
その固定点を持たないliftは
\[
 (x,y)\longmapsto
 \left(\frac{\lambda}{x},
             -\frac{\lambda y}{x^2}\right)
 \tag{Root-2}\label{eq:qa-avatar-root-pair-translation}
\]
であり、$E$上では非零二捩れ点
$(0,0)$によるtranslationである。
この式からも二分割と二捩れ点の
相互復元を直接確認できる。

最後に論文Bの非cubic被約境界は、
二つの二根集合の非順序対を
$\PGL_2$で割った群oidとして記述されており、
四根相異なる開ではまさに第二集合である。
従ってBの分類との合成が最後の
同型類全単射を与える。
これはBのO-marked対象への
逆構成による対応であり、
元のavatar AS代数の直接認識とは区別する。
\end{proof}

\subsection{二根対のquartic resolventと退化avatarの非可逆Picardパラメータ}
\label{subsec:qa-rootpair-resolvent-singular-avatar}

\begin{theorem}[noncubic根対の三重quartic resolventと三つの退化型]
\label{thm:qa-noncubic-rootpair-quartic-resolvent}
$k$を標数零の代数閉体、$W$を二次元ベクトル空間とし、
$\mathcal D_2=\PP(\operatorname{Sym}^2W^\vee)$、
$\mathcal D_4=\PP(\operatorname{Sym}^4W^\vee)$と置く。
二つの二次形式をその順序を忘れて掛ける射
\[
 \mu:\operatorname{Sym}^2(\mathcal D_2)\longrightarrow
             \mathcal D_4,\qquad
 \{[f_1],[f_2]\}\longmapsto[f_1f_2]
 \tag{Res-4}\label{eq:qa-rootpair-resolvent-map}
\]
は全射有限射であり、一般次数は$3$である。
根が相異なる四次形式$q$の上の三点は
四根の非順序二根対への三分割を表す。
アフィンmonic chartで
\[
 q(x)=x^4+a x^3+b x^2+c x+d
\]
とすると、四根$\alpha_1,\ldots,\alpha_4$の
二組分割から得る
$\alpha_i\alpha_j+\alpha_k\alpha_\ell$の
三つの値は
\[
 R_q(z)=z^3-bz^2+(ac-4d)z
             -(a^2d+c^2-4bd)
 \tag{Res-cub}\label{eq:qa-rootpair-resolvent-cubic}
\]
の三根に一致する。
この三重射は古典的なquartic resolventの
幾何的な表示である
\cite[Section 4.3]{Cremona1999}。

$\mathcal D_2^{\mathrm{sf}}$を二つの相異なる
根を持つ二次形式の開部分とする。
論文Bのnoncubic根対表示の
\[
 \mathcal R_{\mathrm{sf}}=
 \operatorname{Sym}^2(\mathcal D_2^{\mathrm{sf}})
\]
上では$q=f_1f_2$の幾何的根配置は
次の三型に限られ、それぞれ
\[
 C_q:\quad y^2=f_1f_2
 \quad\text{in }\operatorname{Tot}\OO_{\PP^1}(2)
 \tag{Res-curve}\label{eq:qa-rootpair-spectral-curve}
\]
の射影的二重被覆が次の形となる。
\[
 \begin{array}{c|c|c}
 \text{根配置}&q\text{の重複度}&C_q\\\hline
 B_1\cap B_2=\varnothing&(1,1,1,1)&\text{smooth genus }1\\
 |B_1\cap B_2|=1&(2,1,1)&\text{既約nodal genus }1\\
 B_1=B_2&(2,2)&\text{二つの}\PP^1\text{が二節点で交わる}
 \end{array}
 \tag{Res-strata}\label{eq:qa-rootpair-three-strata}
\]
三型とも$C_q$は連結なGorenstein曲線で
算術種数$1$である。
この構成は$\operatorname{PGL}(W)$と
根対交換の作用に自然であり、論文Bの
根対スタック
$\mathfrak R=[((\mathcal D_2^{\mathrm{sf}})^2)/
(\operatorname{PGL}(W)\times S_2)]$
からquartic二重被覆のスタックへの射を与える。
ただし全$\operatorname{Sym}^2(\mathcal D_2)$の
有限射と、平方因子を除いた
$\mathcal R_{\mathrm{sf}}$上の制限射は区別する。
\end{theorem}
\begin{proof}
四次形式の任意の因数分解は
その根の重複度付き四点集合から
二点ずつを選ぶ有限個の操作なので、
$\mu$の幾何fiberは有限である。
$\operatorname{Sym}^2(\mathcal D_2)$は射影的、
従って$\mu$はproperかつquasi-finiteで
有限である。代数閉体上ではすべての四次形式が
二次形式の積となるので全射。
四根が相異なれば分割は三つあり、
一般次数は三となる。
式~\eqref{eq:qa-rootpair-resolvent-cubic}は
四根の基本対称式を展開すれば得られる。
例えば三値を
$u_{12|34}=\alpha_1\alpha_2+\alpha_3\alpha_4$
などと置くと、
$\sum u=b$、
$\sum_{r<s}u_ru_s=ac-4d$、
$\prod u=a^2d+c^2-4bd$である。

$f_1,f_2$の根が各々相異なるので、
積の各根の重複度は高々二である。
共通根の個数が$0,1,2$の三場合が
表の各行に一致する。
四単根なら二重被覆はsmoothでRiemann--Hurwitzにより
genus一。一つの二重根の近くでは
$y^2=t^2u(t)$、$u(0)\ne0$なので、
étale局所で$y'^2=t^2$という普通節点になる。
残り二単根から二重被覆は既約で、
正規化はgenus零。
二つの二重根なら$f_1=f_2$（射影類として）で、
$y^2=f_1^2$は二つのsection
$y=\pm f_1$の和である。
各sectionは$\PP^1$で、$f_1$の二単根で横断する。
二重被覆の代数
$\OO_{\PP^1}\oplus\OO_{\PP^1}(-2)y$
から$\chi(\OO_{C_q})=1+(-1)=0$であり、
全型で$p_a(C_q)=1$。
局所的に平面曲線の節点またはsmoothなので
Gorensteinでもある。
群の作用については積と二重被覆の
構成が自然であることから従う。
\end{proof}

\begin{theorem}[根対からの相対torsion-free spectral avatar]
\label{thm:qa-noncubic-rootpair-torsionfree-spectral-module}
論文Bの二形式表示
\[
 (f_1,f_2)\in
 \bigl(\operatorname{Sym}^2W^\vee\bigr)_{\mathrm{nd}}^2
\]
上で、$B=\PP(W)$と置き、
$\widetilde p:\widetilde C\to B$、
$p:C\to B$を以下の有限射、
$r:\widetilde C\to C$をその間の自然な射とする。
\[
 \widetilde C=
 \operatorname{Spec}_{\PP(W)}
  \frac{\OO_{\PP(W)}[c,d]}
       {(c^2-f_1,\ d^2-f_2)},\qquad
 C=\operatorname{Spec}_{\PP(W)}
  \frac{\OO_{\PP(W)}[y]}{(y^2-f_1f_2)}
\]
を取る。ただし$c,d$は局所的に
$\OO_{\PP(W)}(-1)$の生成元、$y=cd$は
$\OO_{\PP(W)}(-2)$の生成元と解釈する。
$\widetilde C$から$C$への有限射は
$(c,d)\mapsto(-c,-d)$の商であり、
$\OO_C$-加群としての反不変部分と
その$B$への直像はそれぞれ
\[
 \mathcal M=(r_*\OO_{\widetilde C})^-,
 \qquad
 p_*\mathcal M
 \cong\OO_B(-1)c\oplus\OO_B(-1)d
 \tag{Res-M}\label{eq:qa-rootpair-spectral-module}
\]
であり、$\mathcal M$は
\[
 y c=f_1 d,\qquad y d=f_2 c
 \tag{Res-MF}\label{eq:qa-rootpair-spectral-matrix-factorization}
\]
を満たす。全二形式族と任意の基底変更にわたり、
$C$は有限平坦二重被覆、$\mathcal M$は
\emph{相対平坦なrank-one torsion-free sheaf}
（各幾何fiber上で次数零）となる。
この構成は二形式の基底変更・スカラー変更・
二形式交換に関して自然なので、
論文Bのnoncubic表示スタック
$\mathfrak N^{\mathrm{pres}}=[P/G]$から
「種数一の有限二重被覆とrank-one
torsion-free sheaf」の相対群oidへの射を与える。

二根対が交わらない滑らかなfiberでは、
$\widetilde C\to C$はétale $\mu_2$-torsor、
$\mathcal M$は非自明な
$\Pic^0(C)[2]$の線束である。
これは定理
~\ref{thm:qa-avatar-two-torsion-root-pair-bijection}
で定めた分割の2-torsionと一致する。
他方、根対が一根を共有するとき
$\mathcal M$はその節点で局所自由でなく、
二根対が一致するときは二つの節点で
局所自由でない。
従って\emph{全noncubic根対配置にわたって}
自然に存在するのは通常の$\Pic^0$線束ではなく、
この相対torsion-free spectral moduleである。
\end{theorem}
\begin{proof}
二つの二次平方根の代数は、
$\OO_B$上の基底$1,c,d,cd$を持つ
階数四の局所自由代数である。
$\widetilde p_*\OO_{\widetilde C}$の
同時符号変換の不変部分は$1,cd$、
反不変部分は$c,d$なので
式~\eqref{eq:qa-rootpair-spectral-module}を得る。
$y=cd$による積は
式~\eqref{eq:qa-rootpair-spectral-matrix-factorization}
で与えられ、形式的には
\[
 \begin{pmatrix}0&f_2\\f_1&0\end{pmatrix}^2
      =f_1f_2 I_2
\]
というmatrix factorizationである。
$p_*\mathcal M$は$B$上で
$\OO(-1)^{\oplus2}$だから、
二形式パラメータとその任意基底変更に対して平坦。
各幾何fiberでは二次拡大のgeneric点でrank一で、
局所的に二次平方根代数の直和因子であるから
torsion-freeである。
Euler特性は
\[
 \chi(C,\mathcal M)
 =\chi(\PP^1,\OO(-1)^{\oplus2})=0
 =\chi(C,\OO_C)
\]
なので次数零である。

$f_1,f_2$に共通零点がなければ、
$c,d$が同時に零となる点がないため
同時符号対合は自由で、
$\widetilde C\to C$は有限étaleの二重torsor。
このとき反不変部分$\mathcal M$は
invertibleかつ$\mathcal M^2\simeq\OO_C$。
$H^0(C,\mathcal M)=H^0(B,p_*\mathcal M)
 =H^0(\PP^1,\OO(-1)^{\oplus2})=0$
なので非自明である。
$\widetilde C$の三つの二次部分被覆の
分岐分割から、これは前定理の2-torsionに一致する。

共有根の局所では単元の平方根をétale局所に取り、
\[
 \OO_{C,o}^{\wedge}
  \simeq k[[x,y]]/(y^2-x^2),\qquad
 \mathcal M_o\cong
   \langle c,d\mid yc=xd,\ yd=xc\rangle .
\]
$\mathcal M_o/(x,y)\mathcal M_o$の次元は二なので、
rank一局所自由ではない。
一致根対では$C$は二つのsectionの和で、
その二つの交点の局所方程式も同じになり、
$\mathcal M$は各節点で非局所自由である。
二形式の交換は$c,d$の交換、
スカラー変更は平方根のスカラー変更として
作用し、構成は元の表示群の作用に整合する。
\end{proof}

\begin{corollary}[noncubic中央二重平方環の可換点曲線からの幾何射]
\label{cor:qa-noncubic-central-square-point-curve-spectral-avatar}
論文A・Bのnoncubic中央二重平方環
\[
 B_f=k[u,v]\langle c,d\rangle/
          (c^2-f_1(u,v),\,d^2-f_2(u,v))
\]
に対し、通常の可換化を施した環と
その射影曲線は
\[
 B_f^{\mathrm{ab}}=
 k[u,v,c,d]/(c^2-f_1,\,d^2-f_2),\qquad
 E_f=\operatorname{Proj}(B_f^{\mathrm{ab}})
 \subset\PP^3
\]
である。これは定理
~\ref{thm:qa-noncubic-rootpair-torsionfree-spectral-module}
の四重被覆$\widetilde C$と自然に同型であり、
\[
 \boxed{\quad
 E_f\xrightarrow{/\,\mu_2}C_{f_1f_2}
 \xrightarrow{\;2:1\;}\PP^1,
 \qquad
 \mathcal M=(r_*\OO_{E_f})^- .
 \quad}
 \tag{NC-point-avatar}\label{eq:qa-noncubic-pointcurve-spectral-factorization}
\]
$\mu_2$は$(c,d)\mapsto(-c,-d)$として作用する。
この対合の固定点は$f_1,f_2$の
共通根の上に限られるので、二根対が
非交差なら第一射はétale、
一根共有または一致根対では
固定点に対応する節点でétaleでなくなる。
この場合、有限平坦な二重torsorと
みなしてはならない。

この射はnoncubic線形代数$B_f$から
rank-one torsion-free spectral avatarへの
\emph{具体的な幾何構成}である。
非可換な$B_f$の$\qgr$と
$\OO$-marked avatar側のheartが同値である
と主張するものではない。
\end{corollary}
\begin{proof}
$u,v$は中央にあるので、可換化において
追加される関係は$cd=dc$だけである。
そのProJの$\PP^1=\operatorname{Proj}k[u,v]$上の
相対座標環は$1,c,d,cd$を基底とする
二重平方根代数となる。
これは前定理の$\widetilde C$そのものである。
同時符号の不変部分は$1,cd$、
反不変部分は$c,d$なので、
商と$\mathcal M$の表示も一致する。
非自明な同時符号の固定点は、
射影座標$(u,v)$が非零である限り
$c=d=0$を満たす点に限られ、
これは$f_1=f_2=0$と同値である。
$(u,v)=(0,0)$では二つの平方関係から
$c=d=0$となり射影点は存在しない。
従って固定点の条件が従う。
\end{proof}

\begin{proposition}[共有根退化での可逆Picard延長の局所障害]
\label{prop:qa-rootpair-nodal-picard-nonextension}
$R=k[[s]]$上の二根対の退化
\[
 f_1=x(x-1),\qquad f_2=(x-s)(x-2)
\]
を取り、$C_R:y^2=f_1f_2$とする。
その一般fiberは滑らかな楕円曲線であり、
特殊fiberは既約な一節点曲線となる。
一般fiberの自然な非零2-torsion線束
$\mathcal M_K$は、
\emph{$C_R$全体の可逆線束には延長できない}。
従って二根対から
$C:y^2=f_1f_2$と
\emph{可逆な}$P\in\Pic^0(C)$を
割り当てる滑らかなavatar辞書は、
この族に沿ってそのまま相対射としては延長できない。

この障害は、$C_R$の節点における
表面$A_1$特異点と、
そこでの非自明なrank-one reflexive module
によって検出される。
$C_R$以外の曲線へのavatar再埋込み、
$\Sigma_2$へのelementary transform、
またはcompactified Picard型データ
\cite{AltmanKleiman1980}による
延長の可能性は排除しない。
ただし同じ一径数基底の上で全空間の$A_1$特異点を
最小解消するだけでは不十分である。
命題~\ref{prop:qa-nodal-resolution-odd-exceptional-degree}の
反不変moduleの可逆化は例外曲線上の次数が奇数であり、
任意の垂直Cartier因子でtwistしても
相対的な多重次数零線束にはならない。
第~\ref{sec:qa-spectral-boundary-models}節の二traitでは，
任意の有限DVR基底変更後の最小正則模型でも二乗自明な線束への延長は存在しない
（定理~\ref{thm:qa-boundary-input-two-torsion-line-obstruction}）．
線束自体への延長と，非可逆TF層・二次積を保持する標的は区別する．
\end{proposition}
\begin{proof}
$s=0$では$f_1f_2=x^2(x-1)(x-2)$であり、
特殊曲線は既約nodal genus一である。
共通根$x=s=0$の完成局所環では、
単元の平方根で座標を正規化すると
\[
 A=k[[x,s,y]]/(y^2-x(x-s)).
\]
二形式を別々に平方根化する被覆の完成環は
$B=k[[c,d]]$で、
\[
 x=c^2,\quad x-s=d^2,\quad y=cd,\quad
 A\simeq B^{\{\pm1\}}
\]
となる。$B$上の同時符号の反不変部分
$M=Ac+Ad$は、商$A$上の
rank-one maximal Cohen--Macaulay、従って
rank-one reflexive moduleである。
その最小生成元は$c,d$の二つなので非自由であり、
$A$は表面$A_1$特異点を持つ。
これは前定理の$\mathcal M$の完成stalkである。

相対曲線の全空間$C_R$はnormalであり、
その唯一の特異点で上記$A_1$局所型を持つ。
$\mathcal M$はnormal曲面上の
rank-one reflexive sheafで、
一般fiberでは$\mathcal M_K\in\Pic^0(C_K)[2]$。
もし可逆線束$L$が全空間上に存在して
$L_K\simeq\mathcal M_K$となれば、
normal曲面上のrank-one reflexive sheafの
Weil因子記述から、$L$と$\mathcal M$の差は
特殊fiberに支持されるWeil因子となる。
この特殊fiberは\emph{既約}かつ
パラメータ$s$で切られたCartier因子なので、
その整数倍は主因子である。
従って$L\simeq\mathcal M$となり、
節点での非局所自由性に矛盾する。
最後の留保は$C_R$を固定したときの
可逆延長障害と、別のavatarモデルの存在を
混同しないためである。
\end{proof}

\begin{theorem}[全noncubic根対と二捩れsquare avatarの体上同型類全単射]
\label{thm:qa-noncubic-rootpair-all-geometric-square-avatar-bijection}
$k$を標数零の代数閉体とする。
論文Bのnoncubic根対の全幾何点、すなわち
$\PP^1$上の\emph{各々は被約}な二点因子の
非順序対$\{B_1,B_2\}$の
$\PGL_2$-同型類集合を考える。
これは、次の条件を持つsquare avatarの
\emph{同型類集合}と自然に全単射する。
\[
 \begin{gathered}
 (Y=\PP^1\times\PP^1;\
  D\in|\OO_Y(2,2)|;\ F_s,G_s,P),\\
 P=G_sF_s^{-1}\ne\OO_D,\qquad
 P^{\otimes2}\simeq\OO_D,\\
 D\text{ はsmooth genus }1,\quad
 \text{既約一節点、または二}\PP^1
 \text{が二節点で交わる曲線}.
 \end{gathered}
 \tag{Sq-all-nc}\label{eq:qa-all-noncubic-square-avatar-bijection}
\]
ここで$F_s=\OO_D(1,0)$、
$G_s=\OO_D(0,1)$であり、
二rulingの順序は保持する。
この全単射は、論文Bの
\cite[\texttt{thm:noncubic-intrinsic-presentation-equivalence}]{QuadricModuli2026}
を介して、\emph{全noncubic O-marked H2幾何同型類}の
集合論的なavatar分類を与える。

構成は次の通りである。
まず$q=f_1f_2$から
式~\eqref{eq:qa-rootpair-spectral-curve}の
二重被覆$p:C_q\to\PP^1$と
$F_s=p^*\OO_{\PP^1}(1)$を取る。
\begin{enumerate}
\item 四根相異なる場合は、
定理~\ref{thm:qa-avatar-two-torsion-root-pair-bijection}
の分割に対応する非零
$P\in\Pic^0(C_q)[2]$を取る。
\item 一根共有または二根対一致の場合は
$\Pic^0(C_q)\simeq\mathbb G_m$なので、
その\emph{唯一の}非自明二捩れ点$P=-1$
を取る。
\end{enumerate}
いずれの場合も$G_s=F_sP$は次数二で
大域生成され、二つの完全線形系の積
\[
 C_q\longrightarrow
 \PP H^0(F_s)^\vee\times\PP H^0(G_s)^\vee
\]
は$(2,2)$反標準Cartier因子への閉埋込みとなる。
この逆操作は、square avatarの
第一rulingの二重被覆$D\to\PP^1$の
分岐四次因子と$P$から、
四根相異なる場合には三分割の一つを、
一根共有・一致根対では唯一の
被約二根対分割を復元することである。

\emph{重要：} これは\emph{体上の同型類の全単射}である。
上の二場合で特別に選んだ
$P=-1\in\mathbb G_m$は、
必ずしも滑らかな族の2-torsion線束の極限でない。
実際、命題
~\ref{prop:qa-rootpair-nodal-picard-nonextension}
の一根共有族では延長できない。
従ってこの全単射を、そのまま
通常のPicard avatarスタックへの
相対モジュライ射、または元のavatar
AS代数の$\qgr$のpointed同値と
読み替えてはならない。
\end{theorem}
\begin{proof}
smoothな場合は定理
~\ref{thm:qa-avatar-two-torsion-root-pair-bijection}
による。以下、特異な二場合を調べる。

\emph{既約一節点の場合。}
$C_q$の正規化は$\nu:\PP^1\to C_q$で、
節点の二逆像を$a,b$とする。
$\Pic^0(C_q)$の線束は、
正規化上の自明線束の
$a,b$での貼合せスカラー
$\lambda\in k^\times$で分類される。
従って$\Pic^0(C_q)\simeq\mathbb G_m$で、
非自明二捩れは貼合せ比$-1$だけである。

$F_s=p^*\OO(1)$は正規化上次数二で、
節点での貼合せ比を$1$に規格化できる。
$G_s=F_sP$は正規化上の$\OO(2)$で、
貼合せ比は$-1$となる。
$H^0(C_q,G_s)$は$H^0(\PP^1,\OO(2))$の
$a,b$における値が符号反対になる
二次元部分空間である。
これはbasepoint-freeである。
実際、$a,b$の二評価汎関数の
非自明な線形結合は、
二次Veroneseの双対円錐上で
単一点の評価汎関数には一致しない。
（異なる二点を結ぶ割線が二次曲線と
交わるのはその両端のみである。）
従って$|G_s|$も有限な次数二射を与える。

$F_s\not\simeq G_s$なので、
この二射の積は総次数二ではあり得ず、
$C_q$上の一般点で単射である。
像の二rulingへの次数は各二だから
$(2,2)$因子で、算術種数は一。
元の$C_q$も算術種数一なので、
有限双有理射の構造層の余核はEuler特性零、
従って零であり、積射は閉埋込みとなる。

\emph{二根対一致の場合。}
$C_q=C_+\cup C_-$は二つの$\PP^1$が
相異なる二節点で交わる曲線である。
正規化上の各成分の次数零線束は自明であり、
二つの節点での貼合せ比
$(\lambda_1,\lambda_2)$を
成分ごとの定数倍で同一視すると、
$\lambda_1/\lambda_2\in\mathbb G_m$
だけが残る。
よって$\Pic^0(C_q)\simeq\mathbb G_m$、
非自明二捩れはこの比が$-1$のものに限る。
$F_s$は各成分で$\OO_{\PP^1}(1)$であり、
$G_s=F_sP$も各成分で次数一、
二節点での値の二条件から
$H^0(C_q,G_s)$は二次元で大域生成される。
二つの成分の像は、
第一$\PP^1$から第二$\PP^1$への
二つの$(1,1)$グラフとなる。
貼合せ比$-1$ゆえ二グラフは異なり、
ちょうど二つの節点で横断する。
従って像は元の$C_q$と同型な
$(2,2)$Cartier因子である。

逆向きは平方freeな二つの二次因子の
分割を数えることで得られる。
四単根なら三通りで非零$E[2]$の三点に対応する。
$(2,1,1)$なら
$(a,b)\mid(a,c)$の一通り、
$(2,2)$なら
$(a,b)\mid(a,b)$の一通りしか
二因子ともsquarefreeにする分割はない。
同じ各曲線の$\Pic^0[2]\setminus0$の
点数が$3,1,1$なので、逆は一意である。
元の第二rulingは
$G_s=F_sP$で復元されるから、
曲線・線束の同型はavatarの同型となる。
最後にBのnoncubic根対分類と合成すれば
O-marked同型類への全単射が得られる。
相対射に関する留保は直前命題の
非延長例により必要である。
\end{proof}

\begin{theorem}[全noncubic O-marked同型類の$\Sigma_2$反標準avatar分類]
\label{thm:qa-all-noncubic-sigma2-anticanonical-avatar-bijection}
$k$を標数零の代数閉体とする。
$\Sigma_2=\PP_{\PP^1}(\OO\oplus\OO(2))$の
負断面を$S=h-2f$とし、
$D\in|2h|=|-K_{\Sigma_2}|$を
$D\cap S=\varnothing$を満たす
相対反標準Cartier因子とする。
このとき、次の二集合は自然に全単射する。
\[
 \boxed{
 \begin{aligned}
 &\{\text{全noncubic O-marked quadricの
             幾何exact同型類}\}\\
 &\quad\longleftrightarrow
 \left\{
 \begin{array}{l}
 (D\subset\Sigma_2,P),\quad
 D\cap S=\varnothing,\\
 0\ne P\in\Pic^0(D)[2]
 \end{array}
 \right\}
/\operatorname{Aut}(\Sigma_2;h,f).
 \end{aligned}}
 \tag{NC-Sigma2}\label{eq:qa-all-noncubic-sigma2-avatar}
\]
右辺で$D$はあらかじめ滑らか、既約、被約と
仮定しない。しかし$P$の非自明二捩れ性から、
実際には$D$はsmooth genus一、
既約一節点、または二つの$\PP^1$が二節点で
交わる型のいずれかになる。
この三型は論文Bの二根対の
非交差／一根共有／一致に正確に対応する。

具体的な同型類の対応は
定理~\ref{thm:qa-noncubic-rootpair-all-geometric-square-avatar-bijection}
のsquare avatar
$(D_s;F_s,G_s,P)$に、
系~\ref{cor:qa-square-chart}の
幾何的変換$\mathsf T^{-1}$を適用したものである。
$\Sigma_2$側の制限線束を
$F=\OO_{\Sigma_2}(f)|_D$、
$H=\OO_{\Sigma_2}(h)|_D$とすると、
\[
 H\simeq F^2,\qquad
 F_s=F,\quad G_s=FP,\quad
 H_s=F_sG_s=F^2P
 \tag{NC-Sigma2-dict}\label{eq:qa-noncubic-sigma2-square-dictionary}
\]
が全対応の明示公式となる。

これは\emph{体上の集合論的全単射}であり、
$\Sigma_2$ avatarから元のAS代数の
pointed heartを直接回復する定理でも、
Picard line bundleを使う相対スタック同値でもない。
smooth域でこのsquare換算を指定AS10 mutationと一番頂点の
pointingへ接続した対象は、
定理~\ref{thm:qa-spectral-avatar-pointed-involution}により
元の対象の二同種双対に対応する。
元とpointed同値となるのは$\chi=0,2$だけであり、
この対応先は二つのnoncubic境界でH2族にも延長できない
（系~\ref{cor:qa-spectral-involution-no-h2-extension}）。
命題~\ref{prop:qa-rootpair-nodal-picard-nonextension}
の共有根退化では、固定されたspectral曲線族上の
可逆$P$への相対延長が不可能なことに注意する。
\end{theorem}
\begin{proof}
$D\cap S=\varnothing$のとき、
$S=h-2f$の定義から
$\OO_D(S)\simeq\OO_D$となり、
$H\simeq F^2$が従う。
また$D$はどの$\PP^1$-fiberも成分として含まない。
実際、fiberは$S$と交わるので、
そのような成分は$D\cap S=\varnothing$に反する。
従って$p:D\to\PP^1$は有限平坦次数二である。
そのtrace-zero直和分解を取ると
\[
 p_*\OO_D\simeq\OO_{\PP^1}\oplus
                \OO_{\PP^1}(-2)
\]
となる。後項の次数は
$\chi(\OO_D)=0$と$p_a(D)=1$から決まる。
よって$D$は四次式$q$を用いた
$y^2=q$型の二重被覆になる。
$q$が四単根、一つの二重根、二つの二重根なら
前定理の三型になる。
それ以外の四次根配置は
$(3,1)$のcuspidal型、
$(4)$の二枝接触型、
または$q=0$の非被約double sectionである。
これらの$\Pic^0$は標数零で
加法型$\mathbb G_a$となり、
非自明二捩れ点を持たない。
この主張は正規化と特異点での単元の貼合せ、
非被約double sectionでは
$1+\OO_{\PP^1}(-2)$の加法的拡大から得られる。
従って$P\ne\OO_D,\ P^2=\OO_D$は
三つのsmooth/nodal/I2型を強制する。

幾何的変換については
系~\ref{cor:qa-square-chart}が、
$\Sigma_2$の負断面を避けるavatarを
square chartへ送る操作$\mathsf T$を
任意基底変更と可換な同値として構成している。
そこで$H=F^2$を代入すると
square側の線束は
\[
 F_s=F,\qquad G_s=HPF^{-1}=FP
\]
となり、$P=G_sF_s^{-1}$かつ$P^2=\OO_D$である。
逆にこの二捩れsquare wallのデータに
$\mathsf T^{-1}$を適用すれば
$H=F_sG_sP^{-1}=F_s^2$となり、
系~\ref{cor:qa-square-chart}と
定理~\ref{thm:qa-open-charts}により
負断面を避ける$\Sigma_2$に入る。
両操作は互いに逆であり、
曲面・ruling・因子・Picard類の同型を保つ。
これと前定理の全noncubic二根対の
幾何同型類全単射、論文Bの
\cite[\texttt{thm:noncubic-intrinsic-presentation-equivalence}]{QuadricModuli2026}
を合成して結論を得る。
相対射に関する留保は直前の
$A_1$局所非延長の結果による。
\end{proof}

\begin{corollary}[noncubic良moduliの同じ点を区別する二つの$\Sigma_2$ avatar]
\label{cor:qa-sigma2-avatar-separates-special-noncubic-classes}
論文Bのnoncubic境界の良moduli
$\chi:\mathfrak N\to\mathbb A^1$において、
一根共有と二根対一致はともに$\chi=1$へ写るが、
これらは異なるpointed exact同型類である。
定理~\ref{thm:qa-all-noncubic-sigma2-anticanonical-avatar-bijection}
の$\Sigma_2$ avatarによれば、
前者は\emph{既約nodal}反標準因子$D$と
唯一の非零二捩れ$P$、
後者は\emph{二$\PP^1$が二節点で交わる}
反標準因子$D$と唯一の非零二捩れ$P$に対応する。
従って曲線$D$の既約成分構成が、
良moduliでは同一視される二つの
O-marked同型類を区別する。
\end{corollary}
\begin{proof}
良moduliの特殊fiberの二幾何同型類と
$\chi=1$への写像は論文B
\cite[\texttt{prop:noncubic-good-moduli-fibers}]{QuadricModuli2026}
による。
根対の共通根の個数に応じた
二重被覆$C_q$の特異型の分類は
定理~\ref{thm:qa-noncubic-rootpair-quartic-resolvent}
にあり、$\Sigma_2$での再埋込みは
前定理で証明した。
既約性・成分数が異なるので
二avatarは同型になり得ない。
\end{proof}

\begin{remark}[節点型avatarの二つの明示$(2,2)$模型]
\label{rem:qa-singular-square-avatar-explicit-models}
定理~\ref{thm:qa-noncubic-rootpair-all-geometric-square-avatar-bijection}
の特異な二型は、いずれも
$\PP^1_X\times\PP^1_Z$上で明示できる。
既約一節点型では、正規化$\PP^1_t$の
$t=1,-1$を貼り合わせ、非自明二捩れの
貼合せ比を$-1$とする。
$F_s$の二切断は$1,t^2$、
$G_s$の二切断は$t,1-t^2$であり、
\[
 X=t^2,\qquad Z=\frac{1-t^2}{t}
 \quad\Longrightarrow\quad
 XZ^2=(1-X)^2 .
 \tag{Sq-node}\label{eq:qa-nodal-square-avatar-model}
\]
これは$(X,Z)=(1,0)$に普通節点を持つ
既約$(2,2)$曲線である。
式は各$\PP^1$で斉次化するものとする。

一致根対型では、$C_+\cup C_-$を二つの
$\PP^1$の和として、その二交点を
$X=0,\infty$に取る。$F_s$は両成分で
$X$座標、非自明二捩れ$P$による
$G_s$は第一成分で$Z=X$、
第二成分で$Z=-X$となる。
従って
\[
 (Z-X)(Z+X)=0
 \tag{Sq-I2}\label{eq:qa-i2-square-avatar-model}
\]
は二つの$(1,1)$曲線が
$(0,0),(\infty,\infty)$で横断する
反標準$(2,2)$因子を与える。
この二模型は、特異な場合の
$P=-1\in\Pic^0(D)[2]$が本当に
square avatarの二rulingを作ることを
直接に示している。
\end{remark}

\begin{proposition}[Picard零切断の全体を古典pointとして追加できない理由]
\label{prop:qa-avatar-zero-picard-collapse}
$Y=\PP^1\times\PP^1$、
$D\in|\OO_Y(2,2)|$を任意の反標準Cartier因子とし、
$P=\OO_D$と置く。
このとき曲線上の四線束
\[
 (\OO_D,\ F_s,\ G_s,\ F_sG_s)
\]
の有向End代数（各積を含む）は、
$D$の選択に依存せず、
$Y$上の可換quadricの標準四項collection
\[
 (\OO_Y,\OO_Y(1,0),\OO_Y(0,1),\OO_Y(1,1))
\]
の有向End代数と同型である。
従って$\Pic^0$の零切断をそのまま
曲線型avatarに追加すると、
異なる$D$が同一の可換foundation代数を与え、
O-marked分類への単射性が失われる。
古典的な可換quadricは、曲線型とは別に
一点として追加する必要がある。
\end{proposition}
\begin{proof}
相異なる四項の間のHomに現れる
$\OO_Y(a,b)$は、
$a,b\in\{-1,0,1\}$の範囲にある。
各々についてCartier制限列
\[
 0\to\OO_Y(a-2,b-2)\to
 \OO_Y(a,b)\to\OO_D(a,b)\to0
\]
を考えると、$a-2\le-1$、$b-2\le-1$であり、
Künnethの公式から
$H^0(Y,\OO(a-2,b-2))=
H^1(Y,\OO(a-2,b-2))=0$である。
従って四項間の全Homの制限は同型となる。
構成は線束の合成に関して自然なので、
有向End代数の積も保存する。
右辺は$D$に依存しない。
\end{proof}

\begin{remark}[二つの $\Sigma _2$ への移行]
\label{rem:qa-rotated-square}
square helix の foundation を一つ進めた rotated square の偏極は
\[
 F_r=F_s,\qquad H_r=F_sG_sP
\]
である（本稿の幾何的証明は定理~\ref{thm:qa-core-avatar-square-rotated}；
原草稿\cite[命題11.1・定理15.1]{AvatarVSeven2026}とも整合する）。
一方、\eqref{eq:qa-square-dictionary}の偏極は $H=F_sG_sP^{-1}$ である。
前者の $\Sigma _2$ 条件は $P=F_sG_s^{-1}$、後者の条件は $P=G_sF_s^{-1}$ となる。
これらは一般には別の壁である。
foundation の移動や mutation とともに選ばれる構造対象も変わるため、
固定された構造対象の内在的一点 Hilbert 関手を同じものとして扱わない。
\end{remark}

\subsection{古典点と最深境界の扱い}

本節の stack は一貫して $P$ が各 fiber で非自明な曲線型を扱った。
体上の分類では、$P=\OO_D$ とした有限 directed algebra は $D$ の選択に依存せず、
\[
 \mathrm{Cl}_0:\operatorname{coh}(\PP^1\times\PP^1),\qquad
 \mathrm{Cl}_2:\operatorname{coh}(\PP(1,1,2))
\]
という二つの古典型に対応する
\cite[注意11.3・定理15.3]{Q12Avatar2026}。
後者の $\PP(1,1,2)$ は重み付き射影 stack である。
これは同型類の説明であり、相対 Picard の零切断を二点へ自動的に縮約する
stack の構成ではない。零切断では $D$ とその自己同型が残るため、
相対的に古典点を加えるには別途の収縮または同一視を定義する必要がある。

最深の十二本矢層についても、三頂点の有限安定表現のモジュライが $\Sigma _2$ になることと、
標準 heart における内在的一点 Hilbert 曲面が $\Sigma _2$ になることは異なる。
この区別のための内在的な検査として，
例~\ref{ex:qa-weighted-not-h2}では重み付き古典模型の全標準頂点で，
例~\ref{ex:qa-hirzebruch-not-h2}では $\coh\Sigma_2$ 自身の全可逆層で，
反Serre軌道の生成条件が破れることを示した。
古典型の負断面上には、一つの foundation では安定であっても、
次の foundation では最後の頂点から生成されなくなる表現が存在する
\cite[命題14.1]{Q12Avatar2026}。
したがって本節の開部分 stack だけから、全 $\mathcal A_{12}^{\times}$ の
内在的 quadric モジュライへの射、あるいはその射の定義域が開であることを結論しない。
そこでは構造対象、標準 heart、豊富性、Serre 平方根および相対 AS 持上げを
別に検証しなければならない。
